% ============================================================
%  يُترجَم بـ XeLaTeX  (في Overleaf: Menu → Compiler → XeLaTeX)
%  الخط: Amiri (متوفر في Overleaf، أو حمّله من github.com/aliftype/amiri)
% ============================================================
\documentclass[12pt,a4paper]{article}
\usepackage[a4paper,margin=1.7cm,top=2.3cm,bottom=2.9cm,headheight=15pt,footskip=1.5cm]{geometry}
\usepackage{amsmath,amssymb}
\usepackage{icomma}
\usepackage[english,bidi=default,layout=lists]{babel}
\babelprovide[import,main]{arabic}
\babelfont{rm}[Scale=1.08]{Amiri}
\babelfont{sf}[Scale=1.08]{Amiri}
\usepackage[most]{tcolorbox}
\usepackage{tikz}
\usetikzlibrary{arrows.meta,calc,shadings}
\usepackage{array}
\usepackage{colortbl}
\usepackage{fancyhdr}
\usepackage{titlesec}
\usepackage{enumitem}
\usepackage{collcell}
\usepackage{needspace}
\usepackage{etoolbox}
\newcolumntype{C}{>{\collectcell\ar}c<{\endcollectcell}}
\pretocmd{\section}{\Needspace{15\baselineskip}}{}{}
\pretocmd{\subsection}{\Needspace{11\baselineskip}}{}{}

\definecolor{cDef}{HTML}{1F4E9C}
\definecolor{cThm}{HTML}{B3261E}
\definecolor{cRem}{HTML}{2E7D32}
\definecolor{cEx}{HTML}{C76B00}
\definecolor{cMeth}{HTML}{6A1B9A}
\definecolor{cSol}{HTML}{00695C}
\definecolor{cSec}{HTML}{1F3A68}

\titleformat{\section}{\Large\bfseries\color{cSec}}{\thesection.}{0.6em}{}[{\color{cSec!60}\titlerule[1pt]}]
\titleformat{\subsection}{\large\bfseries\color{cThm!85!black}}{\thesubsection.}{0.5em}{}
\titlespacing*{\section}{0pt}{14pt}{8pt}
\titlespacing*{\subsection}{0pt}{10pt}{5pt}
\setlength{\parindent}{0pt}
\setlength{\parskip}{4pt}
\setlist{itemsep=2pt,topsep=3pt}

\newtcolorbox{arbox}[3][]{enhanced,
  colback=#2!5,colframe=#2!80!black,arc=2mm,boxrule=0.8pt,
  left=6pt,right=6pt,top=8pt,bottom=5pt,
  fonttitle=\bfseries,coltitle=white,
  title={\foreignlanguage{arabic}{#3}},
  attach boxed title to top right={xshift=-5mm,yshift=-3mm},
  boxed title style={colback=#2!80!black,arc=1.5mm,boxrule=0pt},
  before upper={\selectlanguage{arabic}},
  before skip=8pt,after skip=8pt,breakable,#1}
\newenvironment{BOX}[2]{\begin{otherlanguage}{english}\begin{arbox}{#1}{#2}}{\end{arbox}\end{otherlanguage}}
\newenvironment{tarif}[1][تعريف]{\begin{BOX}{cDef}{#1}}{\end{BOX}}
\newenvironment{khassia}[1][خاصية]{\begin{BOX}{cThm}{#1}}{\end{BOX}}
\newenvironment{molahada}[1][ملاحظة]{\begin{BOX}{cRem}{#1}}{\end{BOX}}
\newenvironment{mithal}[1][مثال]{\begin{BOX}{cEx}{#1}}{\end{BOX}}
\newenvironment{tariqa}[1][طريقة]{\begin{BOX}{cMeth}{#1}}{\end{BOX}}

\newcounter{exo}
\newenvironment{tamrin}[1]{\stepcounter{exo}\begin{BOX}{cSec}{التمرين \theexo: #1}}{\end{BOX}}
\newenvironment{hal}{\Needspace{12\baselineskip}\begin{otherlanguage}{english}\begin{arbox}[breakable]{cSol}{حل التمرين \theexo}}{\end{arbox}\end{otherlanguage}}

\newenvironment{LTR}{\begin{otherlanguage}{english}}{\end{otherlanguage}}
\newcommand{\ar}[1]{\foreignlanguage{arabic}{#1}}

\renewcommand{\dots}{\mathinner{\ldotp\ldotp\ldotp}}
\renewcommand{\ldots}{\mathinner{\ldotp\ldotp\ldotp}}
\renewcommand{\cdots}{\mathinner{\cdotp\cdotp\cdotp}}
\newcommand{\N}{\mathbb{N}}
\newcommand{\Z}{\mathbb{Z}}
\newcommand{\D}{\mathbb{D}}
\newcommand{\Q}{\mathbb{Q}}
\newcommand{\R}{\mathbb{R}}
\newcommand{\pgcd}{\operatorname{PGCD}}
\newcommand{\ppcm}{\operatorname{PPCM}}

% \NLseg{a}{b}{[}{]} ; \NLright{a}{[} ; \NLleft{b}{]}  (interval on a number line)
\newcommand{\NLseg}[4]{\begin{LTR}\begin{tikzpicture}[baseline=-0.6ex]
 \draw[-{Stealth[length=2mm]},gray!70](0,0)--(3.4,0);
 \draw[line width=2pt,cThm!80!black](0.7,0)--(2.7,0);
 \node[cThm!80!black,font=\LARGE] at (0.7,0){$#3$};
 \node[cThm!80!black,font=\LARGE] at (2.7,0){$#4$};
 \node[below=3pt,font=\footnotesize] at (0.7,0){$#1$};
 \node[below=3pt,font=\footnotesize] at (2.7,0){$#2$};
\end{tikzpicture}\end{LTR}}
\newcommand{\NLright}[2]{\begin{LTR}\begin{tikzpicture}[baseline=-0.6ex]
 \draw[-{Stealth[length=2mm]},gray!70](0,0)--(3.4,0);
 \draw[line width=2pt,cThm!80!black](1.3,0)--(3.3,0);
 \node[cThm!80!black,font=\LARGE] at (1.3,0){$#2$};
 \node[below=3pt,font=\footnotesize] at (1.3,0){$#1$};
\end{tikzpicture}\end{LTR}}
\newcommand{\NLleft}[2]{\begin{LTR}\begin{tikzpicture}[baseline=-0.6ex]
 \draw[-{Stealth[length=2mm]},gray!70](0,0)--(3.4,0);
 \draw[line width=2pt,cThm!80!black](0.05,0)--(2.1,0);
 \node[cThm!80!black,font=\LARGE] at (2.1,0){$#2$};
 \node[below=3pt,font=\footnotesize] at (2.1,0){$#1$};
\end{tikzpicture}\end{LTR}}

% ---------- social-media icons (pure TikZ) ----------
\definecolor{cFB}{HTML}{1877F2}
\definecolor{cTG}{HTML}{229ED9}
\definecolor{cIG1}{HTML}{FEDA75}
\definecolor{cIG2}{HTML}{D62976}
\definecolor{cIG3}{HTML}{4F5BD5}
\newcommand{\fbicon}{\begin{LTR}\begin{tikzpicture}[baseline=-0.75ex,x=1em,y=1em]
 \fill[cFB] (0,0) circle (0.55);
 \node[white,inner sep=0,font=\usefont{TU}{lmss}{bx}{n}\small] at (0.04,-0.02) {f};
\end{tikzpicture}\end{LTR}}
\newcommand{\igicon}{\begin{LTR}\begin{tikzpicture}[baseline=-0.75ex,x=1em,y=1em]
 \shade[lower left=cIG1,lower right=cIG2,upper left=cIG2,upper right=cIG3,rounded corners=0.28em] (-0.55,-0.55) rectangle (0.55,0.55);
 \draw[white,line width=0.09em,rounded corners=0.18em] (-0.36,-0.36) rectangle (0.36,0.36);
 \draw[white,line width=0.09em] (0,0) circle (0.17);
 \fill[white] (0.22,0.22) circle (0.05);
\end{tikzpicture}\end{LTR}}
\newcommand{\tgicon}{\begin{LTR}\begin{tikzpicture}[baseline=-0.75ex,x=1em,y=1em]
 \fill[cTG] (0,0) circle (0.55);
 \fill[white] (-0.33,0.0) -- (0.32,0.27) -- (0.2,-0.3) -- (0.0,-0.15) -- (-0.08,-0.28) -- (-0.08,-0.08) -- cycle;
 \fill[cTG!70!black] (-0.08,-0.08) -- (0.25,0.2) -- (0.0,-0.15) -- (-0.08,-0.28) -- cycle;
\end{tikzpicture}\end{LTR}}
% page number: black box, white two-digit number, at the right of the social row
\newcommand{\PageBox}{\begin{LTR}\tikz[baseline=(p.base)]\node[fill=black,text=white,inner sep=2.5pt,minimum width=1.6em,minimum height=1.3em,font=\usefont{TU}{lmss}{bx}{n}\small](p){\ifnum\value{page}<10 0\fi\arabic{page}};\end{LTR}}
% visual left-to-right order: Telegram | Instagram | Facebook | page box
\newcommand{\SocialRow}{\begin{otherlanguage}{english}%
 \ar{\textbf{الأستاذ سوايسية هشام للرياضيات}}\ \tgicon\hfill
 \textbf{prof\_hicham\_math}\ \igicon\hfill
 \ar{\textbf{الأستاذ سوايسية هشام للرياضيات}}\ \fbicon\hspace{2.2em}\PageBox\end{otherlanguage}}

\pagestyle{fancy}
\fancyhf{}
\renewcommand{\headrulewidth}{0.6pt}
\renewcommand{\footrulewidth}{0.4pt}
% \SetHeader{right head: document title}{left head: level -- stream}
\newcommand{\SetHeader}[2]{%
 \fancyhead[C]{\parbox{\textwidth}{\ar{#1\hfill #2}}}%
 \fancyfoot[C]{\parbox{\textwidth}{\small\SocialRow}}}
\usepackage{tkz-tab}
\usepackage{multicol}
% variation / sign table (LTR)
\newenvironment{TAB}{\par\begin{center}\begin{LTR}\begin{tikzpicture}[scale=0.95,every node/.style={font=\small}]}{\end{tikzpicture}\end{LTR}\end{center}}
% grid + axes, clipped: \begin{GR}{xmin}{xmax}{ymin}{ymax}{unit cm}
\newenvironment{GR}[5]{\par\begin{center}\begin{LTR}\begin{tikzpicture}[x=#5cm,y=#5cm,font=\footnotesize]
 \draw[gray!25,very thin] (#1,#3) grid (#2,#4);
 \draw[-{Stealth[length=2mm]}] (#1,0)--(#2+0.3,0) node[below left]{$x$};
 \draw[-{Stealth[length=2mm]}] (0,#3)--(0,#4+0.3) node[below left]{$y$};
 \node[below left] at (0,0){$O$};
 \begin{scope}\clip (#1,#3) rectangle (#2,#4);}
{\end{scope}\end{tikzpicture}\end{LTR}\end{center}}
\newcommand{\Ticks}[2]{\foreach \t in {#1} \node[below,font=\scriptsize] at (\t,0){$\t$}; \foreach \t in {#2} \node[left,font=\scriptsize] at (0,\t){$\t$};}
\newcommand{\bx}[1]{\ensuremath{\colorbox{cDef!12}{$\displaystyle #1$}}}
\newcommand{\TwoCol}[4]{\par\begin{LTR}\noindent\begin{minipage}[t]{#1\textwidth}\selectlanguage{arabic}#2\end{minipage}\hfill\begin{minipage}[t]{#3\textwidth}\selectlanguage{arabic}#4\end{minipage}\end{LTR}\par}
\newcommand{\RL}[4]{\TwoCol{#3}{#4}{#1}{#2}}  % \RL{w1}{right content}{w2}{left content}
\newcommand{\ok}{$\textcolor{cRem}{\checkmark}$}
\newcommand{\limn}{\displaystyle\lim_{n\to+\infty}}
% sub-questions, each on its own line
\newenvironment{sq}{\begin{itemize}[leftmargin=2.4em,itemsep=3pt]}{\end{itemize}}
\newcommand{\qa}{\item[(أ)]}\newcommand{\qb}{\item[(ب)]}\newcommand{\qc}{\item[(جـ)]}\newcommand{\qd}{\item[(د)]}\newcommand{\qe}{\item[(هـ)]}
\newcommand{\pa}{\par\smallskip\textbf{(أ)} }\newcommand{\pb}{\par\smallskip\textbf{(ب)} }\newcommand{\pc}{\par\smallskip\textbf{(جـ)} }\newcommand{\pd}{\par\smallskip\textbf{(د)} }\newcommand{\pe}{\par\smallskip\textbf{(هـ)} }
\newcommand{\QQ}[1]{\par\smallskip\textbf{#1.} }
\newenvironment{tmr}{\stepcounter{exo}\begin{BOX}{cSec}{تمرين محلول \theexo}}{\end{BOX}}
\newcommand{\ABC}[3]{\\[2pt]\begin{tabular}{p{0.28\textwidth}p{0.28\textwidth}p{0.28\textwidth}}(جـ) \ #3 & (ب) \ #2 & (أ) \ #1\end{tabular}}
\pretocmd{\section}{\Needspace{24\baselineskip}}{}{}
\newtcbox{\bxtbox}{on line,colback=cDef!12,colframe=cDef!12,boxrule=0pt,arc=1pt,boxsep=0pt,left=2pt,right=2pt,top=1pt,bottom=1pt}
\newcommand{\bxt}[1]{\foreignlanguage{english}{\bxtbox{\ar{#1}}}}
\usepackage{longtable}
\setlength{\LTpre}{4pt}\setlength{\LTpost}{0pt}
\newcommand{\Phase}[1]{\textbf{#1}}
\newcommand{\Fld}[2]{\textbf{#1:}\ #2}
% رأس المذكرة على نمط النموذج: إطار بعنوان «مذكرة رقم»، معلومات الحصة، ثم إطار الكفاءات
\newcommand{\ModhHead}[9]{%
\begin{LTR}
\begin{tcolorbox}[enhanced,colback=white,colframe=cThm!85!black,arc=4mm,boxrule=1pt,top=14pt,bottom=6pt,left=8pt,right=8pt,
 overlay={\node[fill=cThm!85!black,text=white,rounded corners=3mm,inner xsep=14pt,inner ysep=4pt,font=\large\bfseries] at ([yshift=0pt]frame.north){\ar{مذكرة رقم #1}};},
 before upper={\selectlanguage{arabic}}]
\begin{tabular}{@{}p{0.4\linewidth}p{0.55\linewidth}@{}}
\Fld{المستوى}{#2} & \Fld{الثانوية}{\dotfill}\\[3pt]
\Fld{المادة}{رياضيات} & \Fld{ميدان التعلم}{#3}\\[3pt]
\Fld{المدة}{#5} & \Fld{الوحدة التعلمية}{#4}\\[3pt]
\Fld{الأستاذ}{\dotfill} & \Fld{المحتوى المعرفي}{#6}\\
\end{tabular}
\end{tcolorbox}
\begin{tcolorbox}[enhanced,colback=white,colframe=cThm!85!black,arc=3mm,boxrule=0.8pt,sidebyside,sidebyside align=top,lefthand ratio=0.5,
 segmentation style={dashed,cThm!60},before upper={\selectlanguage{arabic}},before lower={\selectlanguage{arabic}}]
\textbf{الكفاءات المستهدفة:}\par #8
\tcblower
\textbf{الكفاءات القبلية:}\par #7
\end{tcolorbox}
\end{LTR}
\par\smallskip\textbf{الوسائل:} #9\par\smallskip}
\newenvironment{Sayr}{%
\begin{LTR}\renewcommand{\arraystretch}{1.25}
\begin{longtable}{|p{0.17\textwidth}|>{\centering\arraybackslash}p{0.06\textwidth}|p{0.53\textwidth}|>{\centering\arraybackslash}p{0.12\textwidth}|}
\hline\rowcolor{cThm!12}\ar{\textbf{توجيهات وتعاليق}} & \ar{\textbf{المدة}} & \ar{\textbf{سير الحصة}} & \ar{\textbf{مراحل الحصة}}\\\hline\endhead
\hline\endfoot}{\end{longtable}\end{LTR}}
% سطر في الجدول: \Row{مرحلة}{سير}{مدة}{توجيهات}
\newcommand{\Row}[4]{\ar{\small #4} & \ar{#3} & \ar{#2} & \cellcolor{cThm!8}\ar{\Phase{#1}}\\\hline}

\newcommand{\Lab}[2]{\foreignlanguage{english}{\tcbox[on line,colback=#1,colframe=#1,boxrule=0pt,arc=2pt,boxsep=1pt,left=5pt,right=5pt,top=1pt,bottom=1pt]{\ar{\textcolor{white}{\textbf{#2}}}}}}
\newcommand{\Num}[1]{\foreignlanguage{english}{\tikz[baseline=(c.base)]\node[fill=cDef,text=white,inner sep=2pt,font=\footnotesize\bfseries](c){#1};}}
\newcommand{\ck}{\ensuremath{{\color{cRem}\checkmark}}}
\newenvironment{Kad}[1]{\par\smallskip\Lab{#1}}{\par}
\SetHeader{\small\bfseries\color{cSec}مذكرات محور: الدالتان الأسية واللوغاريتمية}{\small\color{cSec}السنة الثالثة ثانوي -- علوم تجريبية}
\begin{document}
\begin{LTR}
\begin{tcolorbox}[enhanced,colback=cThm!85!black,colframe=cThm!85!black,arc=3mm,halign=center,boxrule=0pt,top=12pt,bottom=12pt,before upper={\selectlanguage{arabic}}]
{\color{white}\fontsize{24}{30}\selectfont\bfseries مذكرات محور: الدالتان الأسية واللوغاريتمية}\\[8pt]
{\color{white}\large السنة الثالثة ثانوي -- علوم تجريبية (صالحة لشعبتي رياضيات وتقني رياضي)}\\[4pt]
{\color{white}\large 8 مذكرات (16 ساعة) -- السنة الدراسية 2026-2027}
\end{tcolorbox}
\end{LTR}
\begin{BOX}{cThm}{يتضمن هذا العمل}
\begin{enumerate}[label=\textbf{مذكرة \arabic{enumi}:},leftmargin=6em,start=10]
\item الدالة الأسية النيبيرية: تعريف وخواص جبرية.
\item دراسة الدالة الأسية ، معادلات ومتراجحات.
\item الدوال $x\mapsto e^{kx}$ والنمذجة.
\item دراسة الدالة $x\mapsto e^{u(x)}$.
\item الدالة اللوغاريتمية النيبيرية: تعريف وخواص جبرية.
\item دراسة الدالة $\ln$ ، معادلات ومتراجحات.
\item الدالة $x\mapsto\ln u(x)$ ، اللوغاريتم العشري.
\item المعادلات التفاضلية $y^{\prime}=ay+b$.
\end{enumerate}
مراجع النشاطات والتمارين: الكتاب المدرسي للسنة الثالثة ثانوي ، الجزء الأول (باب الدوال الأسية واللوغاريتمية: نشاطا ص 75 و 77 ، تمارين ص 102--112). مذكرات الدالة الأسية (10--13) لا تستعمل $\ln$.
\end{BOX}
\newpage
\ModhHead{10}{3 ع ت}{تحليل}{الدالتان الأسية واللوغاريتمية}{02 سا}{الدالة الأسية النيبيرية: تعريف وخواص جبرية}%
{\begin{itemize}\item الاشتقاقية ، المعادلات التفاضلية $y^{\prime}=f$.\end{itemize}}%
{\begin{itemize}\item التعرف على الدالة الأسية وخواصها الجبرية.\end{itemize}}%
{المنهاج ، الكتاب المدرسي ، مراجع ، الحاسبة البيانية. مجدول.}
\begin{Sayr}
\Row{مرحلة الانطلاق}{%
\textbf{التهيئة النفسية:} هل توجد دالة تساوي مشتقتها؟}{10 د}{سؤال تحفيزي.}
\Row{التشخيص والاكتشاف}{%
\Lab{cEx}{مناقشة نشاط 01 صفحة 75}\par\ck\ نقبل وجود دالة $f$ حيث $f^{\prime}=f$ و $f(0)=1$ ، ونقرب قيمها بطريقة أولر: $f(x+h)\approx(1+h)f(x)$.\par
\ck\ نبرهن أنها لا تنعدم ، وأنها وحيدة ، وأن $f(a+b)=f(a)f(b)$.\par
}{35 د}{طريقة أولر بالمجدول.}
\Row{بناء المعارف}{%
\Lab{cDef}{تعريف}\ \textbf{الدالة الأسية} $\exp$ هي الدالة الوحيدة القابلة للاشتقاق على $\mathbb R$ حيث $\exp^{\prime}=\exp$ و $\exp(0)=1$ ؛ نكتب $\exp(x)=e^x$ مع $e=\exp(1)\approx2{,}718$.\par\Lab{cThm}{خاصية}\ من أجل كل $a$ و $b$ حقيقيين و $n$ صحيح:\par\ck\ $e^{a+b}=e^ae^b$ ؛ $e^{-a}=\dfrac1{e^a}$ ؛ $e^{a-b}=\dfrac{e^a}{e^b}$ ؛ $(e^a)^n=e^{na}$.\par
\ck\ $e^x>0$ من أجل كل $x$.\par
}{50 د}{الخواص تُحفظ.}
\Row{}{%
\Lab{cEx}{مثال}\ $\dfrac{e^{2x}\times e^{-x+1}}{e^x}=e^{2x-x+1-x}=e$.\par\Lab{cThm!80!black}{حذار}\ $e^{a+b}\neq e^a+e^b$.}{}{}
\Row{التقويم}{%
\Lab{cEx}{تطبيق}\par بسّط: $A=(e^x)^2\times e^{-2x+1}$ ؛ $B=\dfrac{e^{3x}-e^x}{e^x}$.\par\textbf{الحل:} $A=e$ ؛ $B=e^{2x}-1$.\par\textbf{واجب منزلي:} تمرينا 2 و 3 صفحة 102.\par\smallskip\textbf{ملاحظات حول سير الحصة:} \dotfill}{25 د}{تقويم فردي.}
\end{Sayr}
\newpage
\ModhHead{11}{3 ع ت}{تحليل}{الدالتان الأسية واللوغاريتمية}{02 سا}{دراسة الدالة الأسية ، معادلات ومتراجحات}%
{\begin{itemize}\item الدالة الأسية وخواصها.\end{itemize}}%
{\begin{itemize}\item دراسة الدالة الأسية وتمثيلها ، وحل معادلات ومتراجحات أسية.\end{itemize}}%
{المنهاج ، الكتاب المدرسي ، مراجع ، الحاسبة البيانية.}
\begin{Sayr}
\Row{مرحلة الانطلاق}{%
\textbf{التهيئة النفسية:} ما إشارة مشتقة $e^x$ ؟ ماذا تستنتج؟}{10 د}{سؤال.}
\Row{التشخيص والاكتشاف}{%
\Lab{cEx}{وضعية مقترحة}\par\ck\ $\exp^{\prime}=\exp>0$: متزايدة تماماً ؛ جدول قيم بالحاسبة لـ $x=\pm10$ ، $\pm20$.\par
\ck\ تخمين النهايتين عند $+\infty$ و $-\infty$.\par
}{35 د}{الحاسبة البيانية.}
\Row{بناء المعارف}{%
\Lab{cThm}{خاصية}\par\ck\ $\exp$ متزايدة تماماً على $\mathbb R$ ؛ $\lim\limits_{x\to+\infty}e^x=+\infty$ ، $\lim\limits_{x\to-\infty}e^x=0$ (محور الفواصل مقارب).\par
\ck\ المماس عند $0$: $y=x+1$.\par
\ck\ $e^a=e^b\Leftrightarrow a=b$ ؛ $e^a<e^b\Leftrightarrow a<b$.\par
\Lab{cEx}{مثال}\par\ck\ $e^{2x-1}=e^{x+3}$: $2x-1=x+3$ ، $x=4$.\par
\ck\ $e^{x^2}\geqslant e^4$: $x^2\geqslant4$ ، $x\in]-\infty;-2]\cup[2;+\infty[$.\par
}{50 د}{جدول التغيرات والمنحني.}
\Row{}{%
\Lab{cThm!80!black}{حذار}\ المعادلة $e^x=k$ مع $k\leqslant0$ لا حل لها.}{}{}
\Row{التقويم}{%
\Lab{cEx}{تطبيق}\par حل في $\mathbb R$: $e^{2x}-(e+1)e^x+e=0$ (ضع $X=e^x$).\par\textbf{الحل:} $X^2-(e+1)X+e=0$: $X=1$ أو $X=e$ ، أي $x=0$ أو $x=1$.\par\textbf{واجب منزلي:} تمارين من باب الدالة الأسية.\par\smallskip\textbf{ملاحظات حول سير الحصة:} \dotfill}{25 د}{تقويم فردي.}
\end{Sayr}
\newpage
\ModhHead{12}{3 ع ت}{تحليل}{الدالتان الأسية واللوغاريتمية}{02 سا}{الدوال $x\mapsto e^{kx}$ والنمذجة}%
{\begin{itemize}\item دراسة الدالة الأسية.\end{itemize}}%
{\begin{itemize}\item دراسة الدوال $x\mapsto e^{kx}$ واستعمالها في نمذجة ظواهر.\end{itemize}}%
{المنهاج ، الكتاب المدرسي ، مراجع ، الحاسبة البيانية.}
\begin{Sayr}
\Row{مرحلة الانطلاق}{%
\textbf{التهيئة النفسية:} كمية دواء في الدم تتناقص بسرعة متناسبة مع الكمية الموجودة: كيف نعبر عنها؟}{10 د}{وضعية طبية.}
\Row{التشخيص والاكتشاف}{%
\Lab{cEx}{وضعية مقترحة}\par\ck\ ظاهرة سرعة تغيرها متناسبة مع قيمتها: $N^{\prime}=kN$.\par
\ck\ الدالة $N(t)=N_0e^{kt}$ تحقق ذلك.\par
}{35 د}{نمذجة.}
\Row{بناء المعارف}{%
\Lab{cThm}{خاصية}\par\ck\ $(e^{kx})^{\prime}=k\,e^{kx}$.\par
\ck\ $k>0$: متزايدة ، $\lim\limits_{x\to+\infty}e^{kx}=+\infty$ ؛ $k<0$: متناقصة ، $\lim\limits_{x\to+\infty}e^{kx}=0$.\par
\Lab{cEx}{مثال}\ $f(x)=e^{-2x}$: $f^{\prime}(x)=-2e^{-2x}<0$ ، متناقصة ؛ $N(t)=N_0e^{-0{,}2t}$ تتناقص نحو $0$.\par}{50 د}{مقارنة منحنيات حسب $k$.}
\Row{}{%
\Lab{cThm!80!black}{حذار}\ $e^{kx}\neq k\,e^x$.}{}{}
\Row{التقويم}{%
\Lab{cEx}{تطبيق}\par $g(x)=3e^{0{,}5x}$. ادرس اتجاه تغيرها ، واكتب معادلة مماس منحناها عند $0$.\par\textbf{الحل:} $g^{\prime}(x)=1{,}5e^{0{,}5x}>0$: متزايدة ؛ $g(0)=3$ ، $g^{\prime}(0)=1{,}5$: $y=1{,}5x+3$.\par\textbf{واجب منزلي:} تمارين 11 ، 12 ، 13 ، 15 صفحة 112.\par\smallskip\textbf{ملاحظات حول سير الحصة:} \dotfill}{25 د}{تقويم فردي.}
\end{Sayr}
\newpage
\ModhHead{13}{3 ع ت}{تحليل}{الدالتان الأسية واللوغاريتمية}{02 سا}{دراسة الدالة $x\mapsto e^{u(x)}$}%
{\begin{itemize}\item مشتق دالة مركبة ، الدالة الأسية.\end{itemize}}%
{\begin{itemize}\item حساب مشتق $e^{u}$ ودراسة اتجاه تغيرها.\end{itemize}}%
{المنهاج ، الكتاب المدرسي ، مراجع ، الحاسبة البيانية.}
\begin{Sayr}
\Row{مرحلة الانطلاق}{%
\textbf{التهيئة النفسية:} كيف نشتق $e^{x^2}$ ؟}{10 د}{سؤال.}
\Row{التشخيص والاكتشاف}{%
\Lab{cEx}{وضعية مقترحة}\par\ck\ $e^{u}=\exp\circ u$: مشتق دالة مركبة: $u^{\prime}\times\exp(u)$.\par
}{35 د}{تطبيق القاعدة.}
\Row{بناء المعارف}{%
\Lab{cThm}{خاصية}\ إذا كانت $u$ قابلة للاشتقاق على $I$: $(e^{u})^{\prime}=u^{\prime}e^{u}$ ، و $e^u$ لها نفس اتجاه تغير $u$.\par\Lab{cEx}{مثال}\par\ck\ $f(x)=e^{x^2-2x}$: $f^{\prime}(x)=(2x-2)e^{x^2-2x}$: متناقصة ثم متزايدة ، قيمة حدية صغرى $f(1)=e^{-1}$.\par
\ck\ $g(x)=(x-1)e^x$: $g^{\prime}(x)=xe^x$ ، قيمة حدية صغرى $g(0)=-1$.\par
}{50 د}{جدول التغيرات.}
\Row{}{%
\Lab{cThm!80!black}{حذار}\ إشارة $u^{\prime}e^u$ هي إشارة $u^{\prime}$ لأن $e^u>0$.}{}{}
\Row{التقويم}{%
\Lab{cEx}{تطبيق}\par ادرس اتجاه تغير $h(x)=e^{\frac1x}$ على $]0;+\infty[$.\par\textbf{الحل:} $h^{\prime}(x)=-\dfrac1{x^2}e^{\frac1x}<0$: متناقصة.\par\textbf{واجب منزلي:} تمارين من باب الدالة الأسية.\par\smallskip\textbf{ملاحظات حول سير الحصة:} \dotfill}{25 د}{تقويم فردي.}
\end{Sayr}
\newpage
\ModhHead{14}{3 ع ت}{تحليل}{الدالتان الأسية واللوغاريتمية}{02 سا}{الدالة اللوغاريتمية النيبيرية: تعريف وخواص جبرية}%
{\begin{itemize}\item الدالة الأسية ، مبرهنة القيم المتوسطة.\end{itemize}}%
{\begin{itemize}\item تعريف الدالة $\ln$ واستعمال خواصها الجبرية.\end{itemize}}%
{المنهاج ، الكتاب المدرسي ، مراجع ، الحاسبة البيانية.}
\begin{Sayr}
\Row{مرحلة الانطلاق}{%
\textbf{التهيئة النفسية:} ما العدد $x$ الذي يحقق $e^x=2$ ؟}{10 د}{سؤال مفتوح.}
\Row{التشخيص والاكتشاف}{%
\Lab{cEx}{مناقشة نشاط 02 صفحة 77}\par\ck\ $\exp$ مستمرة ومتزايدة تماماً من $0$ إلى $+\infty$: لكل $y>0$ عدد وحيد $x$ حيث $e^x=y$.\par
\ck\ نسميه $\ln y$ ؛ $\ln2\approx0{,}693$.\par
}{35 د}{نشاط الكتاب.}
\Row{بناء المعارف}{%
\Lab{cDef}{تعريف}\ من أجل $x>0$ و $y$ حقيقي: $y=\ln x\Leftrightarrow e^y=x$.\par\Lab{cThm}{خاصية}\par\ck\ $\ln1=0$ ، $\ln e=1$ ، $e^{\ln x}=x$ ($x>0$) ، $\ln(e^x)=x$.\par
\ck\ $\ln(ab)=\ln a+\ln b$ ؛ $\ln\dfrac ab=\ln a-\ln b$ ؛ $\ln a^n=n\ln a$ ؛ $\ln\sqrt a=\frac12\ln a$ ($a,b>0$).\par
}{50 د}{الخواص تُحفظ مع أمثلة مضادة.}
\Row{}{%
\Lab{cEx}{مثال}\ $\ln8-2\ln2=3\ln2-2\ln2=\ln2$ ؛ $e^x=5\Leftrightarrow x=\ln5$.\par\Lab{cThm!80!black}{حذار}\ $\ln(a+b)\neq\ln a+\ln b$ ، و $\ln$ معرفة على $]0;+\infty[$ فقط.}{}{}
\Row{التقويم}{%
\Lab{cEx}{تطبيق}\par احسب: $A=\ln(e^2)+\ln\dfrac1e$ ؛ $B=\ln12-\ln3-2\ln2$.\par\textbf{الحل:} $A=2-1=1$ ؛ $B=\ln\dfrac{12}{3\times4}=0$.\par\textbf{واجب منزلي:} تمرين 60 صفحة 106.\par\smallskip\textbf{ملاحظات حول سير الحصة:} \dotfill}{25 د}{تقويم فردي.}
\end{Sayr}
\newpage
\ModhHead{15}{3 ع ت}{تحليل}{الدالتان الأسية واللوغاريتمية}{02 سا}{دراسة الدالة $\ln$ ، معادلات ومتراجحات}%
{\begin{itemize}\item خواص $\ln$ ، دراسة الدالة الأسية.\end{itemize}}%
{\begin{itemize}\item دراسة الدالة $\ln$ وتمثيلها ، وحل معادلات ومتراجحات لوغاريتمية.\end{itemize}}%
{المنهاج ، الكتاب المدرسي ، مراجع ، الحاسبة البيانية.}
\begin{Sayr}
\Row{مرحلة الانطلاق}{%
\textbf{التهيئة النفسية:} منحنيا $e^x$ و $\ln x$ بالحاسبة البيانية: ما العلاقة بينهما؟}{10 د}{رسم بالحاسبة.}
\Row{التشخيص والاكتشاف}{%
\Lab{cEx}{وضعية مقترحة}\par\ck\ $M(a;b)$ على منحني $\exp$ يعني $M_1(b;a)$ على منحني $\ln$: تناظر بالنسبة إلى $y=x$.\par
}{35 د}{التناظر.}
\Row{بناء المعارف}{%
\Lab{cThm}{خاصية}\par\ck\ $\ln^{\prime}(x)=\dfrac1x$: متزايدة تماماً على $]0;+\infty[$.\par
\ck\ $\lim\limits_{x\to0^+}\ln x=-\infty$ (محور التراتيب مقارب) ، $\lim\limits_{x\to+\infty}\ln x=+\infty$ ؛ المماس عند $1$: $y=x-1$.\par
\ck\ $\ln a=\ln b\Leftrightarrow a=b$ ؛ $\ln a<\ln b\Leftrightarrow a<b$ ($a,b>0$).\par
\par\smallskip\begin{LTR}\centering \begin{tikzpicture}[x=0.75cm,y=0.75cm,font=\scriptsize]
\draw[->] (-3.5,0)--(4,0) node[right]{$x$}; \draw[->] (0,-3)--(0,4.3) node[above]{$y$};
\draw[cRem,dashed] (-3,-3)--(4,4) node[right]{$y=x$};
\draw[cDef,thick,domain=-3.5:1.4,samples=60] plot (\x,{exp(\x)}); \node[cDef] at (1.6,4.2){$y=e^x$};
\draw[cThm,thick,domain=0.05:4,samples=80] plot (\x,{ln(\x)}); \node[cThm] at (3.8,1.7){$y=\ln x$};
\draw[gray,thin] (-1.5,-0.5)--(1.5,2.5); \fill (0,1) circle (1.4pt); \fill (1,0) circle (1.4pt);
\end{tikzpicture}\end{LTR}\par \Lab{cEx}{مثال}\par\ck\ $\ln(2x-1)=\ln(x+2)$ ، مع $x>\frac12$: $x=3$.\par
\ck\ $\ln(x-1)<0$: $0<x-1<1$ ، أي $1<x<2$.\par
}{50 د}{مجموعة التعريف أولاً.}
\Row{}{%
\Lab{cThm!80!black}{حذار}\ نعين أولاً مجموعة تعريف المعادلة.}{}{}
\Row{التقويم}{%
\Lab{cEx}{تطبيق}\par حل: $\ln x+\ln(x-2)=\ln3$.\par\textbf{الحل:} الشرط $x>2$ ؛ $x(x-2)=3$: $x=3$ أو $x=-1$ (مرفوض) ؛ الحل $x=3$.\par\textbf{واجب منزلي:} تمرينا 67 و 68 صفحة 107.\par\smallskip\textbf{ملاحظات حول سير الحصة:} \dotfill}{25 د}{تقويم فردي.}
\end{Sayr}
\newpage
\ModhHead{16}{3 ع ت}{تحليل}{الدالتان الأسية واللوغاريتمية}{02 سا}{الدالة $x\mapsto\ln u(x)$ ، اللوغاريتم العشري}%
{\begin{itemize}\item دراسة $\ln$ ، مشتق دالة مركبة.\end{itemize}}%
{\begin{itemize}\item دراسة الدالة $\ln u$ ، وتعريف اللوغاريتم العشري.\end{itemize}}%
{المنهاج ، الكتاب المدرسي ، مراجع ، الحاسبة البيانية.}
\begin{Sayr}
\Row{مرحلة الانطلاق}{%
\textbf{التهيئة النفسية:} الـ pH محلول: $\mathrm{pH}=-\log[\mathrm{H_3O^+}]$ ، ما $\log$ ؟}{10 د}{وضعية من الكيمياء.}
\Row{التشخيص والاكتشاف}{%
\Lab{cEx}{وضعية مقترحة}\par\ck\ $f(x)=\ln(x^2+1)$ معرفة على $\mathbb R$ لأن $x^2+1>0$.\par
\ck\ مشتقها بقاعدة الدالة المركبة: $\dfrac{2x}{x^2+1}$.\par
}{35 د}{تطبيق القاعدة.}
\Row{بناء المعارف}{%
\Lab{cThm}{خاصية}\ إذا كانت $u$ موجبة تماماً وقابلة للاشتقاق على $I$: $(\ln u)^{\prime}=\dfrac{u^{\prime}}u$ ، و $\ln u$ لها نفس اتجاه تغير $u$.\par\Lab{cDef}{اللوغاريتم العشري}\ $\log x=\dfrac{\ln x}{\ln10}$ ($x>0$) ؛ $\log10^n=n$.\par}{50 د}{جدول ملخص.}
\Row{}{%
\Lab{cEx}{مثال}\ $\log1000=3$ ؛ $[\mathrm{H_3O^+}]=10^{-4}$: $\mathrm{pH}=4$.\par\Lab{cThm!80!black}{حذار}\ شرط $\ln u$: $u(x)>0$.}{}{}
\Row{التقويم}{%
\Lab{cEx}{تطبيق}\par ادرس اتجاه تغير $g(x)=\ln(2x-4)$ على $]2;+\infty[$.\par\textbf{الحل:} $g^{\prime}(x)=\dfrac2{2x-4}=\dfrac1{x-2}>0$: متزايدة.\par\textbf{واجب منزلي:} تمارين من باب الدوال اللوغاريتمية.\par\smallskip\textbf{ملاحظات حول سير الحصة:} \dotfill}{25 د}{تقويم فردي.}
\end{Sayr}
\newpage
\ModhHead{17}{3 ع ت}{تحليل}{الدالتان الأسية واللوغاريتمية}{02 سا}{المعادلات التفاضلية $y^{\prime}=ay+b$}%
{\begin{itemize}\item الدوال $e^{kx}$ ، المعادلات $y^{\prime}=f$.\end{itemize}}%
{\begin{itemize}\item حل معادلات تفاضلية من الشكل $y^{\prime}=ay+b$ مع شرط ابتدائي.\end{itemize}}%
{المنهاج ، الكتاب المدرسي ، مراجع ، الحاسبة البيانية.}
\begin{Sayr}
\Row{مرحلة الانطلاق}{%
\textbf{التهيئة النفسية:} قانون تبريد كوب قهوة: سرعة تغير درجة الحرارة متناسبة مع الفرق مع درجة الغرفة. كيف نكتب ذلك؟}{10 د}{وضعية فيزيائية.}
\Row{التشخيص والاكتشاف}{%
\Lab{cEx}{وضعية مقترحة}\par\ck\ $y^{\prime}=2y$: الدوال $Ce^{2x}$ حلول.\par
\ck\ $y^{\prime}=-y+4$: الدالة الثابتة $4$ حل خاص ، ونضيف حلول $y^{\prime}=-y$.\par
}{35 د}{حل خاص + حل المتجانسة.}
\Row{بناء المعارف}{%
\Lab{cThm}{خاصية}\par\ck\ حلول $y^{\prime}=ay$ ($a\neq0$): $y=Ce^{ax}$.\par
\ck\ حلول $y^{\prime}=ay+b$: $y=Ce^{ax}-\dfrac ba$ ، ويوجد حل وحيد يحقق $y(x_0)=y_0$.\par
\Lab{cEx}{مثال}\ $y^{\prime}=-y+4$ و $y(0)=1$: $y=Ce^{-x}+4$ ، $C=-3$: $y=4-3e^{-x}$.\par}{50 د}{التحقق بالتعويض.}
\Row{}{%
\Lab{cThm!80!black}{حذار}\ نتحقق من الحل بالتعويض في المعادلة.}{}{}
\Row{التقويم}{%
\Lab{cEx}{تطبيق}\par حل $y^{\prime}=0{,}5y-1$ مع $y(0)=0$.\par\textbf{الحل:} $y=Ce^{0{,}5x}+2$ ؛ $C=-2$: $y=2-2e^{0{,}5x}$.\par\textbf{واجب منزلي:} تمارين 102 ، 105 ، 106 صفحتي 109 و 110.\par\smallskip\textbf{ملاحظات حول سير الحصة:} \dotfill}{25 د}{نهاية المحور: فرض أو مراجعة.}
\end{Sayr}
\end{document}
